\documentclass[10pt]{article}

\usepackage[utf8]{inputenc}

\usepackage[T1]{fontenc}

\usepackage{amsmath}

\usepackage{amsfonts}

\usepackage{amssymb}

\usepackage[version=4]{mhchem}

\usepackage{stmaryrd}

\usepackage{bbold}



\begin{document}

$\omega(\sigma, f)$ - модуль непрерывности фуункции $f(t)$ на отрезке $[0,1]$, то для частных сумм $S_{n}(t, f)$ порядка $n$ ряда Фурье-Хаара функции $f(t)$ справедливо неравенство



$\sup _{0 \leqslant t \leqslant 1}\left|f(t)-S_{n}(t, f)\right| \leqslant 12 \omega(1 / n, f), \quad n=1,2, \ldots$



X. с. является базисом в пространстве $L_{p}[0,1], 1 \leqslant$ $\leqslant p<\infty$. Если $f(t) \in L_{p}[0,1]$ и $\omega_{p}(\delta, f)$ - интегральный модуль непрерывности функции $f(t)$ в метрике пространства $L_{p}[0,1]$, то (cM. [3])



$\left\|f(t)-S_{n}(t, f)\right\|_{L_{p}[0,1]} \leqslant 24 \omega_{p}(1 / n, f), n=1,2, \ldots$. $\mathrm{X}$. с. является безусловным базисом в пространстве $L_{p}\lceil 0,1]$ при $1<p<\infty$.



Если функция $f(t)$ интегрируема по Лебегу на отрезке $[0,1]$, то ее ряд Фурье-Хаара сходится к ней в любой точке Лебега этой функции и, в частности, почти всюду на $[0,1]$. При әтом сходимость (и абсолютная сходимость) ряда Фурье-Хаара в фиксированной точке отрезка $[0,1]$ зависит липь от значений функции в любой сколь угодно малой окрестности этой тотки.



Для рядов Фурье - Хаара существенно отличаются друг от друга следующие свойства: а) абсолютная сходимость всюду; б) абсолютная сходимость почти всюду; в) абсолютная сходимость на множестве полољительной меры; г) абсолютная сходимость ряда коәффициентов Фурье. Для тригонометрич. рядов все әти свойства равносильны.



Свойства коәффициентов Фурье - Хаара резко отличаются от свойств тригонометрич. коэффициентов Фурье. Напр., если функция $f(t)$ непрерывна на отрезке $[0,1]$, а $a_{n}(f)$ - ее коэффициенты Фурье по системе $\left\{\chi_{n}(t)\right\}$, то справедливо неравенство



\[

\left|a_{n}(f)\right| \leqslant(1 / \sqrt{2 n}) \omega(1 / n, f), n \geqslant 2,

\]



откуда следует, что



\[

a_{n}(f)=o\left(n^{-1 / 2}\right), n \longrightarrow \infty \text {. }

\]



В то же время коэффициенты Фурье - Хаара непрерывных функций не могут убывать слипгко быстро если функция $f(t)$ непрерывна на отрезке $[0,1]$ и



\[

a_{n}(f)=o\left(n^{-3 / 2}\right)

\]



то $f(t) \equiv$ const на $[0,1]$.



Для функций $f(t) \in L^{p}[0,1], 1 \leqslant p<\infty$, справедливы следугщие оценки (см. [3]):



\[

\begin{aligned}

&\left|a_{n}(f)\right| \leqslant n^{(1 / p)-(1 / 2)} \omega_{p}(1 / n, f), n=1,3, \ldots, \\

&\left(\sum_{k=2^{n}+1}^{2^{n+1}}\left|a_{n}(f)\right| p\right)^{1 / p} \leqslant \\

& \leqslant 8 \cdot 2^{n(1 / p-1 / 2)} \omega_{p}\left(1 / 2^{n}, f\right), \quad n=0,1, \ldots .

\end{aligned}

\]



Если же $f(t)$ имеет конечную вариацию $V(f)$ на отрезке $[0,1]$, то



\[

\sum_{k=2^{n+1}}^{2^{n+1}}\left|a_{n}(f)\right| \leqslant\left(3 / 2 \sqrt{2^{n}}\right) V(f), n=0,1, \ldots

\]



Все эти неравенства являются точными в смысле порядка убывания их правых частей при $n \rightarrow \infty$ (в соответствующих классах) (см. [3]).



Интересной особенностью отличаются ряды вида



\[

\sum_{n=1}^{\infty} a_{n} \chi_{n}(t)

\]



безусловно сходящиеся почти всюду: если ряд вида (\textit{) при любом порядке следования его членов сходится почти всюду на множестве $E \subset[0,1]$ положительной меры Лебега (исключительное множество меры нуль может зависеть от порядка следования членов ряда (})), то этот ряд сходится абсолютно ночті всюду на $[0,1]$. Для рядов вида (\textit{) справедлив следующий критерий: чтобы ряд (}) сходился почти всюду на измеримом множестве $E \subset[0,1]$, необходимо и достаточно, чтобы ряд $\sum_{n=1}^{\infty} a_{n}^{2} \chi_{n}^{2}(t)$ сходился потти всюду на $E$.



Ряды Хаара могут служить для представления измеримых функций: для любой конечной почти всюду на отрезке $[0,1]$ измеримой функции $f(t)$ существует ряд вида $(*)$, к-рый почти всюду на $[0,1]$ сходится к функции $f(t)$. При этом конечность функции $f(t)$ существенна: не существует ряда вида $(*)$, схоцящегося к $+\infty($ или $-\infty)$ на множестве положительной меры Лебега.



Jum.: [1] H a a r A., "Math. Ann.», 1910, Bd 69, S. 331-71; [2] А л е к с и ч Г, Проблемы сходимости ортогональных рядов, 1964, т. 63, № 3, c. 356--91, [4] е г о Ж е, там Же, 1967, т. 72, ,$N 2$, с. $193-225,[5]$ Г о л у б о в Б. И., в сб Итоги науки Математический анализ 1970 , М., 1971, с $109-43$



ХААРА УСЛОВИЕ - условие на непрерывные линейно независимые на ограниченном замкнутом множестве $M$ евклидова пространства функции $x_{k}(t), k=1, \ldots, n$. Сформулировано А. Хааром ([1]). Х. у. гарантирует лля любой непрерывной на $M$ функции $f(t)$ единственность по ли но ма на и лу г ші е г о пр и б лиже н і я (н. п.) по системе $\left\{x_{k}(t)\right\}$, т. е. полинома



\[

P_{n-1}(t)=\sum_{k=1}^{n} c_{k} x_{k}(t)

\]



для кі-рого



\[

\begin{gathered}

\max _{t \in M}\left|f(t)-P_{n-1}(t)\right|= \\

=\min _{\left\{a_{k}\right\}} \max _{t \in M}\left|f(t)-\sum_{k=1}^{n} a_{k} x_{k}(t)\right| .

\end{gathered}

\]



X. у. состоит в том, что любой ветривиальный полином вида (*) должен иметь в $M$ не более $n-1$ различных нулей. Чтобы для любой нешрерывной на $M$ функции $f(t)$ существовал единственный полином н. П. по системе $\left\{x_{k}(t)\right\}_{k=1}^{n}$, необходимо и достаточно, чтобы эта система удовлетворяла Х. у. Систему функций, удовлетворяющих Х.у., наз. Чебышева системой. Для таких систем справедлива Чебышева теорема и Валле Пуссена теорема (об альтернансе). Х. у. достаточно для единственности полинома н. П. по системе $\left\{x_{k}(t)\right\}_{k=1}^{n}$ в метрике $L[a, b](M=[a, b])$ для любой непрерывной на $[a, b]$ функции.



Jum. [1] H a a r A , «Math. Ann », 1918, Bd 78, [2] A x и eе р Н. И., Јекции по теории аппроксимации, 2 изд., М., 1965.



ХАДВИГЕРА ГИПОТЕЗА - задача комбинаторной геометрии о покрытии выпуклого тела фигурами специального вида, выдвинутая X. Хадвигером [1]. Пусть $K$ - выпуклое тело $n$-мерного евклидова пространства $\mathbb{R}^{n}$, а $b(K)$ - минимальное тисло тел, гомотетичных $K$ с коәфффициентом гомотетии $k, 0<k<1$, достатотное для покрытия тела $K$. Х. г. заключается в следующем: для любого ограниченного $K \subset \mathbb{R}^{n}$ справедливы неравенства



\[

n+1 \leqslant b(K) \leqslant 2^{n} .

\]



Причем неравенство $b(K)=2^{n}$ характеризует параллелепипед (см. [1]). Х. г. подтверждена для случая $n \leqslant 2$; для $n \geqslant 3$ имеются (1984) лишь частные результаты. Напр., для любого $n$-мерного ограниченного многогранника $K \subset \mathbb{R} n$, каждые две вершины к-рого принадлежат двум различным параллельным опорным гиперплоскостям к $K$, справедливы неравенства (*). Причем $b(K)$ совпадает с числом вершин $K$, а в множестве таких многогранников равенство $b(K)=2^{n}$ проверяется только для параллслепипеда. Этот результат связан с решением одной из Эрдёша задац о числе точек в $\mathbb{R} n$, каждые три из к-рых образуют не тупоугольный треугольник. X. г. связана и с покрытия, разбиения и освещения задачами. Напр., Х. г. может быть рассмотрена как обобщение Борсука проблемь





\end{document}